\section{Hardware-Aware Machine Learning}

A growing body of work recognises that generic efficiency metrics such
as FLOP count or parameter count do not reliably predict real-world
inference latency~\cite{devlin2019bert}. Factors including memory
bandwidth, SIMD utilisation, and cache behaviour can cause a model with
fewer FLOPs to run slower than a model with more FLOPs on a given
hardware platform.

Hardware-aware neural architecture search (NAS) addresses this by
incorporating direct latency measurements or hardware-specific cost
models into the search objective. Platforms such as ProxylessNAS and
MnasNet use reinforcement learning or differentiable search to explore
architecture spaces while constraining latency on a specific target
device.
